\chapter{Estado del Arte}

El análisis computacional de \textit{Whole Slide Images} (WSI) en patología digital ha pasado, en pocos años, desde enfoques especializados para tareas puntuales hacia un paradigma centrado en modelos fundacionales. Este cambio responde a dos características estructurales del dominio. En primer lugar, las WSI son imágenes de resolución extremadamente alta, usualmente del orden de gigapíxeles, lo que impide procesarlas de manera directa y obliga a trabajar mediante parches, agregación jerárquica y estrategias de representación multi-escala. En segundo lugar, la anotación experta en patología es costosa, lenta y altamente dependiente del contexto clínico, por lo que los enfoques que logran aprender representaciones útiles a partir de grandes volúmenes de datos no etiquetados han adquirido una relevancia central \parencite{Basak2023}.

En este escenario, los modelos fundacionales han emergido como una vía particularmente prometedora, ya que permiten preentrenar \textit{encoders} visuales o multimodales sobre grandes colecciones de WSI y luego reutilizar esas representaciones en tareas diversas, tales como clasificación de tejido, subtipificación tumoral, predicción pronóstica, recuperación de casos similares, segmentación y generación de descripciones clínicas. A diferencia de los enfoques tradicionales, donde cada tarea exigía diseñar y entrenar un modelo específico, los modelos fundacionales buscan aprender una representación general del dominio histopatológico que pueda transferirse a múltiples contextos con una cantidad limitada de supervisión adicional. Esta promesa de generalización ha convertido a la patología digital en uno de los escenarios más activos para el desarrollo de modelos fundacionales en medicina.

Sin embargo, el interés por estos modelos no se explica solo por su rendimiento. También se relaciona con la posibilidad de abordar uno de los principales cuellos de botella del área: la heterogeneidad de los datos. Las WSI exhiben variabilidad asociada a protocolos de fijación y tinción, escáneres, aumentos, instituciones, poblaciones de pacientes y criterios diagnósticos. En consecuencia, cualquier representación fundacional que aspire a ser útil en práctica clínica debe ser robusta frente a esta diversidad. Por ello, el verdadero aporte del estado del arte no se agota en identificar qué modelo obtiene mejores resultados agregados, sino en comprender qué tan generales son sus representaciones, bajo qué supuestos fueron entrenadas y qué limitaciones persisten cuando se los evalúa en escenarios realistas de desplazamiento de distribución, subgrupos poblacionales y cambios de adquisición.

\section{Modelos fundacionales en patología digital}

\subsection{MI-Zero}

\textcite{lu2023visuallanguagepretrainedmultiple} proponen MI-Zero, que combina preentrenamiento visión-lenguaje con aprendizaje de múltiples instancias (MIL) para operar sobre \textit{slides} completos y realizar clasificación \textit{zero-shot} sobre WSI gigapíxel, formulando categorías diagnósticas mediante descripciones textuales. Su aporte es trasladar al dominio histopatológico la alineación imagen-texto como supervisión semántica débil transferible a nuevas tareas.

Desde el punto de vista conceptual, MI-Zero muestra que la alineación entre imagen y texto puede utilizarse para clasificación \textit{zero-shot} sobre WSI gigapíxel: en lugar de requerir conjuntos de entrenamiento etiquetados para cada tarea, el paradigma clásico de la patología computacional, explora la posibilidad de formular categorías diagnósticas mediante descripciones textuales y proyectarlas sobre las representaciones visuales del \textit{slide}. Su limitación relevante para esta tesis es que no profundiza en la composición de los datos que sostienen ese comportamiento: la generalización funcional no garantiza homogeneidad entre instituciones, subgrupos o condiciones de adquisición.

\subsection{UNI}

\textcite{chen2024uni} presentan UNI, un \textit{encoder} visual auto-supervisado entrenado exclusivamente sobre imágenes de patología, que demuestra que el dominio dispone de la escala suficiente para sostener un modelo fundacional y que la representación aprendida se transfiere a un amplio conjunto de tareas diagnósticas sin rediseñar el \textit{backbone} en cada caso. Su aporte es consolidar el preentrenamiento auto-supervisado a gran escala como infraestructura común de la patología computacional.

UNI consolida la idea de que el aprendizaje auto-supervisado a gran escala puede capturar regularidades morfológicas suficientemente generales como para ser útiles en tareas diversas: en la práctica, desplaza el centro de gravedad del desarrollo desde la ingeniería de arquitecturas para tareas individuales hacia la calidad del preentrenamiento y la riqueza del conjunto de datos, sugiriendo que la morfología tisular contiene estructura reutilizable entre órganos y objetivos diagnósticos diferentes. Para esta tesis, su relevancia radica en esa reutilización entre dominios, pero la representatividad de las poblaciones y los posibles sesgos de sus datos de preentrenamiento permanecen fuera de su alcance.


\subsection{CONCH}

Entre los modelos desarrollados específicamente para vincular la información morfológica de las imágenes histopatológicas con descripciones semánticas de carácter clínico destaca CONCH (\textit{CONtrastive Learning from Captions for Histopathology}). Presentado por \textcite{Lu2024}, este enfoque combina aprendizaje de representaciones visuales y modelado multimodal mediante el uso de pares imagen-texto del dominio histopatológico, preentrenado sobre más de 1,17 millones de pares mediante aprendizaje contrastivo. Su principal contribución no radica únicamente en mejorar el desempeño en tareas específicas, sino en incorporar una estrategia de supervisión más informativa, donde el lenguaje funciona como un mecanismo para relacionar patrones tisulares con conceptos diagnósticos y contexto clínico relevante.

Su importancia radica en dos puntos. Por una parte, demuestra que los modelos visión-lenguaje pueden entrenarse de forma efectiva en patología con beneficios concretos para clasificación, segmentación, recuperación y descripción semántica. Por otro lado, establece que la patología digital no debe entenderse únicamente como un problema visual, ya que gran parte del conocimiento experto del dominio se expresa en reportes, descripciones diagnósticas, nomenclaturas y criterios anatomopatológicos. Además, si el texto utilizado durante el preentrenamiento proviene de reportes clínicos reales, la representación hereda las regularidades y sesgos del lenguaje clínico institucional (nomenclatura, estilos de redacción, prevalencias), extendiendo el análisis de \textit{bias} más allá de la imagen; esta es la limitación relevante para esta tesis.

\subsection{BiomedGPT}

\textcite{Zhang2024} presentan BiomedGPT, el primer modelo fundacional de visión y lenguaje de código abierto para biomedicina, que integra imágenes, lenguaje y tareas heterogéneas en una misma arquitectura y supera al estado del arte en 16 de 25 experimentos. Su aporte es demostrar la convergencia hacia sistemas biomédicos generalistas capaces de operar entre disciplinas.

En el ámbito de patología digital, esta tendencia es importante por dos razones. Primero, porque las WSI rara vez se interpretan de manera aislada en el flujo clínico real: la lectura anatomopatológica se conecta con antecedentes clínicos, descripciones macroscópicas, biomarcadores, decisiones terapéuticas y reportes textuales. Segundo, porque la utilidad futura de los modelos fundacionales en patología probablemente no se limitará a extraer \textit{features}, sino que incluirá funciones de apoyo a la interpretación, síntesis de hallazgos, búsqueda de casos análogos e interacción con conocimiento biomédico estructurado. Su limitación relevante para esta tesis es que la amplitud de sus datos de entrenamiento dificulta determinar el origen de sus decisiones y evaluar la representatividad de las poblaciones incluidas, reforzando la necesidad de estrategias de trazabilidad y auditoría de datos.

\subsection{Virchow}

\textcite{vorontsov2024virchow} presentan Virchow, un modelo fundacional específicamente diseñado para patología computacional. Parte de la idea de que las \textit{whole slide images} contienen una enorme diversidad de patrones morfológicos que difícilmente pueden capturarse con modelos entrenados en colecciones pequeñas o poco variadas; para enfrentar este problema, los autores emplean aprendizaje auto-supervisado sobre una colección de WSI de gran escala, utilizando una arquitectura \textit{Vision Transformer} entrenada con DINOv2, con el objetivo de aprender representaciones visuales generales que luego puedan transferirse a múltiples tareas posteriores.

La relevancia de Virchow no se limita a su desempeño empírico: sus representaciones se evalúan en detección de cáncer, predicción de biomarcadores y \textit{benchmarks} a nivel de parche, mostrando que una representación aprendida de forma auto-supervisada puede adaptarse eficazmente a problemas clínicamente relevantes sin requerir rediseñar el extractor visual para cada caso, reforzando el desplazamiento desde modelos específicos para tareas aisladas hacia representaciones fundacionales reutilizables. Para esta tesis, su limitación es análoga a la de los modelos anteriores: su evaluación no aborda la heterogeneidad de las condiciones de adquisición ni la composición demográfica de las cohortes.

\section{De la capacidad fundacional a la evaluación rigurosa}

\subsection{FairMedFM}

Parte importante de la literatura ya no se concentra solo en proponer nuevos modelos, sino en evaluar sistemáticamente cómo se comportan los modelos fundacionales en presencia de disparidades entre subgrupos. En esta línea, \textcite{NEURIPS2024_c9826b9e} presentan FairMedFM, un benchmark de \textit{fairness} para modelos fundacionales de imágenes médicas. Aunque su alcance es multimodalidad médica en sentido amplio y no exclusivamente patología digital, demuestra que las métricas agregadas pueden ocultar diferencias sustantivas de rendimiento entre grupos definidos por atributos sensibles o factores clínicamente relevantes.

El valor de FairMedFM radica en que introduce una evaluación estandarizada donde la utilidad global del modelo y la equidad por subgrupos se analizan de manera conjunta. Esta idea es especialmente importante para patología digital, donde el énfasis histórico ha estado puesto en AUC, exactitud o desempeño promedio en tareas específicas, mientras que las brechas entre instituciones, cohortes o grupos demográficos suelen quedar en segundo plano. FairMedFM muestra que esta omisión no es menor: modelos aparentemente sólidos pueden presentar comportamientos desiguales al cambiar el grupo de análisis, lo cual tiene implicancias clínicas y éticas directas.

\subsection{Documentación de conjuntos de datos y trazabilidad del sesgo}

Un aspecto cada vez más visible en la literatura reciente es que la auditoría de \textit{bias} depende críticamente de la calidad de la documentación de los conjuntos de datos. Si un \textit{dataset} no reporta con suficiente detalle variables demográficas, centro de origen, protocolo de tinción, escáner, criterios de inclusión y exclusión, o distribución de clases por subgrupo, entonces la evaluación de \textit{fairness} queda severamente limitada.

Esta línea de investigación evidencia que el avance hacia modelos fundacionales de gran escala no solo plantea desafíos algorítmicos, sino también cuestiones relacionadas con la transparencia de los datos que sustentan su entrenamiento. La amplitud y heterogeneidad de los conjuntos de datos utilizados favorecen la capacidad de transferencia de las representaciones aprendidas, pero al mismo tiempo dificultan la identificación de posibles sesgos, desequilibrios o limitaciones de cobertura. Por ello, la caracterización rigurosa de los conjuntos de datos adquiere un papel central para comprender e interpretar adecuadamente el comportamiento de estos modelos.

Desde esta perspectiva, el estado del arte permite identificar una brecha muy concreta. Existen numerosos trabajos dedicados a desarrollar mejores representaciones visuales o multimodales para WSI, pero son comparativamente menos los que describen con la misma profundidad las propiedades de los datos que alimentan esas representaciones. Esta asimetría resulta especialmente relevante porque sugiere que el desempeño de los modelos no depende únicamente de la arquitectura empleada o de la estrategia de entrenamiento, sino también de características subyacentes de los conjuntos de datos que con frecuencia permanecen insuficientemente documentadas. En consecuencia, aspectos como la cobertura, la representatividad y los posibles sesgos de los datos adquieren un papel central para comprender adecuadamente los resultados reportados.

\section{Brecha de investigación}

Finalmente, aunque los modelos fundacionales han mostrado resultados prometedores en patología digital, persiste una brecha importante en el análisis sistemático de los conjuntos de datos de \textit{whole slide images} que sustentan su entrenamiento y evaluación. En particular, aún falta una caracterización estructurada de variables demográficas, condiciones de adquisición y desbalances de representación, así como una evaluación clara de cómo estos factores pueden afectar la generalización y la equidad de los modelos. En este contexto, esta tesis busca abordar dicha brecha mediante un enfoque centrado en los datos, orientado a identificar, describir y analizar posibles fuentes de \textit{bias} en conjuntos de datos de WSI de patología digital, y a estudiar su relación con el desempeño de modelos fundacionales representativos. De esta forma, el aporte del trabajo no radica en proponer un nuevo modelo, sino en contribuir con una metodología de análisis que permita comprender mejor cómo la composición de los conjuntos de datos influye en el desarrollo de modelos más robustos, confiables y equitativos.
